% Insert with % Insert with % Insert with % Insert with \input{rational-partition-algorithm} before \appendix.
% Uses only the amsart preamble of Diagonal Rational Partitions.
\section{Exact computation by denominator elimination}
\label{sec:exact-computation}

Fix \(N\ge2\) and an integer \(M\ge0\). We compute all counts
\(R_N(m)\), \(0\le m\le M\), in one pass. The last entry for \(N=M=n\)
is \(R(n)\). If the remaining denominators divide \(\Lambda\), their
parts have total in \(\Lambda^{-1}\mathbb Z\). A partial sum that can
be completed to an integer must therefore belong to the same lattice.

Let \(b_1,\ldots,b_d\) be a permutation of \(2,\ldots,N\), where
\(d=N-1\), and put
\[
  \Lambda_d=1,\qquad
  \Lambda_{i-1}=\operatorname{lcm}(b_i,\Lambda_i)
  \quad (i=d,d-1,\ldots,1).
\]
Thus \(\Lambda_i\) is the least common multiple of the denominators
following \(b_i\). We order the denominators by decreasing greatest
prime factor, so that after the group for \(p\) all remaining
denominators are coprime to \(p\).

\subsection{The recurrence}
Let \(f_j\) count the retained configurations of total \(j/D\).
To process denominator \(b\), refine the lattice to denominator
\(H=\operatorname{lcm}(D,b)\). Each allowed part \(a/b\) now has
integer weight \(s=aH/b\). After all parts with denominator \(b_i\)
have been included, the partial sums eligible for retention lie in
\begin{equation}
  H^{-1}\mathbb Z\cap\Lambda_i^{-1}\mathbb Z
  =\gcd(H,\Lambda_i)^{-1}\mathbb Z.
  \label{eq:exact-lattice-intersection}
\end{equation}
The projection therefore keeps every \(H/\gcd(H,\Lambda_i)\)-th
coefficient, starting at index zero.

\smallskip
\noindent\begin{minipage}{\linewidth}
\textsc{RationalPartitionRow}\((N,M)\)
\begin{quote}
\small
\begin{tabbing}
\hspace{1.2em}\=\hspace{1.2em}\=\hspace{1.2em}\=\kill
Choose \(b_1,\ldots,b_d\) and compute \(\Lambda_d,\ldots,\Lambda_0\).\\
\(D\gets1\), \quad \(f\gets(1,0,\ldots,0)\) of length \(M+1\).\\
\textbf{for} \(i=1,\ldots,d\) \textbf{do}\\
\> \(b\gets b_i\), \quad \(H\gets\operatorname{lcm}(D,b)\).\\
\> Set \(g_0,\ldots,g_{MH}\) to zero.\\
\> \textbf{for} \(j=0,\ldots,MD\) \textbf{do} \(g_{jH/D}\gets f_j\).\\
\> \textbf{for} \(1\le a\le\min(N,Mb)\) with \(\gcd(a,b)=1\) \textbf{do}\\
\>\> \(s\gets aH/b\).\\
\>\> \textbf{for} \(j=s,\ldots,MH\) in increasing order \textbf{do}\\
\>\>\> \(g_j\gets g_j+g_{j-s}\).\\
\> \(D\gets\gcd(H,\Lambda_i)\), \quad \(h\gets H/D\).\\
\> \(f\gets(g_{kh})_{k=0}^{MD}\).\\
\textbf{return} \(f\).
\end{tabbing}
\end{quote}
\end{minipage}

\begin{proposition}\label{prop:exact-projection}
The algorithm returns \((R_N(m))_{m=0}^{M}\) for every ordering of the
denominators.
\end{proposition}
\begin{proof}
The initial array counts the empty multiset, and refinement carries
each total \(j/D\) to index \(jH/D\). For a new part of weight \(s\),
let \(u_j,v_j\) be the counts before and after its inclusion. Removing
one copy from each multiset that contains the new part gives
\[
  v_j=u_j+v_{j-s},\qquad v_k=0\quad(k<0).
\]
The increasing loop implements this recurrence. Each multiset is
counted once, since every part is introduced once.

Fix a partition of \(m\le M\), with subtotal \(x\) after \(b_i\).
Since \(m-x\in\Lambda_i^{-1}\mathbb Z\) and \(m\) is an integer,
\(x\in\Lambda_i^{-1}\mathbb Z\). If its previous subtotal was retained,
the updates reach \(x\) on the lattice \(H^{-1}\mathbb Z\).
By \eqref{eq:exact-lattice-intersection}, \(Hx\) is a multiple of
\(h\), so this subtotal is retained too. Positivity gives \(x\le m\),
which justifies truncation at \(M\). By induction every partition of
\(m\) survives. Finally \(\Lambda_d=1\), hence \(D=1\) and \(f_m=R_N(m)\).
\end{proof}

\subsection{The count at order three}
Take the denominators in the order \(3,2\). The thirds must have an
integer subtotal, since \(\gcd(3,2)=1\). To obtain a subtotal \(j\)
from \(1/3,2/3\), their multiplicities satisfy \(u+2v=3j\). For
\(1/2,3/2\) the corresponding equation is \(u+3v=2j\). Thus the two
counts are \(T_j=\lfloor3j/2\rfloor+1\) and
\(B_{2j}\), where \(B_s=\lfloor s/3\rfloor+1\) counts a subtotal
\(s/2\) from the half parts. For \(j=0,1,2,3\) this gives
\[
 (T_j)=(1,2,4,5),\qquad (B_{2j})=(1,1,2,3).
\]
Pairing complementary subtotals gives
\[
 R(3)=\sum_{j=0}^{3}T_jB_{2(3-j)}
     =1\cdot3+2\cdot2+4\cdot1+5\cdot1=16.
\]

\subsection{The count at order four}
Use the order \(3,4,2\). The thirds again have an integer subtotal,
since the remaining denominators divide \(4\). There is now a third
part with denominator \(3\), namely \(4/3\). Its inclusion gives the
count
\[
 U_j=[z^{3j}]\frac{1}{(1-z)(1-z^2)(1-z^4)}.
\]
For the quarter parts \(1/4,3/4\), the next denominator is \(2\),
so \eqref{eq:exact-lattice-intersection} retains half-integer
subtotals. For instance, two copies of \(1/4\) can be completed to
\(1\) by one copy of \(1/2\). The number of ways to make \(r/2\)
from \(1/4,3/4\) is
\[
 Q_r=\#\{(u,v)\in\mathbb Z_{\ge0}^2:u+3v=2r\}
     =\left\lfloor\frac{2r}{3}\right\rfloor+1.
\]
The half parts are again \(1/2,3/2\). Thus the combined count for
integer subtotal \(j\) is \(V_j=\sum_{r=0}^{2j}Q_r B_{2j-r}\), giving
\[
 \begin{array}{c|rrrrr}
 j   &0&1&2&3&4\\ \hline
 U_j &1&2&6&9&16\\
 V_j &1&4&12&27&48
 \end{array}
\]
and the final convolution is
\[
 R(4)=\sum_{j=0}^{4}U_jV_{4-j}
     =1\cdot48+2\cdot27+6\cdot12+9\cdot4+16\cdot1=226.
\]
 before \appendix.
% Uses only the amsart preamble of Diagonal Rational Partitions.
\section{Exact computation by denominator elimination}
\label{sec:exact-computation}

Fix \(N\ge2\) and an integer \(M\ge0\). We compute all counts
\(R_N(m)\), \(0\le m\le M\), in one pass. The last entry for \(N=M=n\)
is \(R(n)\). If the remaining denominators divide \(\Lambda\), their
parts have total in \(\Lambda^{-1}\mathbb Z\). A partial sum that can
be completed to an integer must therefore belong to the same lattice.

Let \(b_1,\ldots,b_d\) be a permutation of \(2,\ldots,N\), where
\(d=N-1\), and put
\[
  \Lambda_d=1,\qquad
  \Lambda_{i-1}=\operatorname{lcm}(b_i,\Lambda_i)
  \quad (i=d,d-1,\ldots,1).
\]
Thus \(\Lambda_i\) is the least common multiple of the denominators
following \(b_i\). We order the denominators by decreasing greatest
prime factor, so that after the group for \(p\) all remaining
denominators are coprime to \(p\).

\subsection{The recurrence}
Let \(f_j\) count the retained configurations of total \(j/D\).
To process denominator \(b\), refine the lattice to denominator
\(H=\operatorname{lcm}(D,b)\). Each allowed part \(a/b\) now has
integer weight \(s=aH/b\). After all parts with denominator \(b_i\)
have been included, the partial sums eligible for retention lie in
\begin{equation}
  H^{-1}\mathbb Z\cap\Lambda_i^{-1}\mathbb Z
  =\gcd(H,\Lambda_i)^{-1}\mathbb Z.
  \label{eq:exact-lattice-intersection}
\end{equation}
The projection therefore keeps every \(H/\gcd(H,\Lambda_i)\)-th
coefficient, starting at index zero.

\smallskip
\noindent\begin{minipage}{\linewidth}
\textsc{RationalPartitionRow}\((N,M)\)
\begin{quote}
\small
\begin{tabbing}
\hspace{1.2em}\=\hspace{1.2em}\=\hspace{1.2em}\=\kill
Choose \(b_1,\ldots,b_d\) and compute \(\Lambda_d,\ldots,\Lambda_0\).\\
\(D\gets1\), \quad \(f\gets(1,0,\ldots,0)\) of length \(M+1\).\\
\textbf{for} \(i=1,\ldots,d\) \textbf{do}\\
\> \(b\gets b_i\), \quad \(H\gets\operatorname{lcm}(D,b)\).\\
\> Set \(g_0,\ldots,g_{MH}\) to zero.\\
\> \textbf{for} \(j=0,\ldots,MD\) \textbf{do} \(g_{jH/D}\gets f_j\).\\
\> \textbf{for} \(1\le a\le\min(N,Mb)\) with \(\gcd(a,b)=1\) \textbf{do}\\
\>\> \(s\gets aH/b\).\\
\>\> \textbf{for} \(j=s,\ldots,MH\) in increasing order \textbf{do}\\
\>\>\> \(g_j\gets g_j+g_{j-s}\).\\
\> \(D\gets\gcd(H,\Lambda_i)\), \quad \(h\gets H/D\).\\
\> \(f\gets(g_{kh})_{k=0}^{MD}\).\\
\textbf{return} \(f\).
\end{tabbing}
\end{quote}
\end{minipage}

\begin{proposition}\label{prop:exact-projection}
The algorithm returns \((R_N(m))_{m=0}^{M}\) for every ordering of the
denominators.
\end{proposition}
\begin{proof}
The initial array counts the empty multiset, and refinement carries
each total \(j/D\) to index \(jH/D\). For a new part of weight \(s\),
let \(u_j,v_j\) be the counts before and after its inclusion. Removing
one copy from each multiset that contains the new part gives
\[
  v_j=u_j+v_{j-s},\qquad v_k=0\quad(k<0).
\]
The increasing loop implements this recurrence. Each multiset is
counted once, since every part is introduced once.

Fix a partition of \(m\le M\), with subtotal \(x\) after \(b_i\).
Since \(m-x\in\Lambda_i^{-1}\mathbb Z\) and \(m\) is an integer,
\(x\in\Lambda_i^{-1}\mathbb Z\). If its previous subtotal was retained,
the updates reach \(x\) on the lattice \(H^{-1}\mathbb Z\).
By \eqref{eq:exact-lattice-intersection}, \(Hx\) is a multiple of
\(h\), so this subtotal is retained too. Positivity gives \(x\le m\),
which justifies truncation at \(M\). By induction every partition of
\(m\) survives. Finally \(\Lambda_d=1\), hence \(D=1\) and \(f_m=R_N(m)\).
\end{proof}

\subsection{The count at order three}
Take the denominators in the order \(3,2\). The thirds must have an
integer subtotal, since \(\gcd(3,2)=1\). To obtain a subtotal \(j\)
from \(1/3,2/3\), their multiplicities satisfy \(u+2v=3j\). For
\(1/2,3/2\) the corresponding equation is \(u+3v=2j\). Thus the two
counts are \(T_j=\lfloor3j/2\rfloor+1\) and
\(B_{2j}\), where \(B_s=\lfloor s/3\rfloor+1\) counts a subtotal
\(s/2\) from the half parts. For \(j=0,1,2,3\) this gives
\[
 (T_j)=(1,2,4,5),\qquad (B_{2j})=(1,1,2,3).
\]
Pairing complementary subtotals gives
\[
 R(3)=\sum_{j=0}^{3}T_jB_{2(3-j)}
     =1\cdot3+2\cdot2+4\cdot1+5\cdot1=16.
\]

\subsection{The count at order four}
Use the order \(3,4,2\). The thirds again have an integer subtotal,
since the remaining denominators divide \(4\). There is now a third
part with denominator \(3\), namely \(4/3\). Its inclusion gives the
count
\[
 U_j=[z^{3j}]\frac{1}{(1-z)(1-z^2)(1-z^4)}.
\]
For the quarter parts \(1/4,3/4\), the next denominator is \(2\),
so \eqref{eq:exact-lattice-intersection} retains half-integer
subtotals. For instance, two copies of \(1/4\) can be completed to
\(1\) by one copy of \(1/2\). The number of ways to make \(r/2\)
from \(1/4,3/4\) is
\[
 Q_r=\#\{(u,v)\in\mathbb Z_{\ge0}^2:u+3v=2r\}
     =\left\lfloor\frac{2r}{3}\right\rfloor+1.
\]
The half parts are again \(1/2,3/2\). Thus the combined count for
integer subtotal \(j\) is \(V_j=\sum_{r=0}^{2j}Q_r B_{2j-r}\), giving
\[
 \begin{array}{c|rrrrr}
 j   &0&1&2&3&4\\ \hline
 U_j &1&2&6&9&16\\
 V_j &1&4&12&27&48
 \end{array}
\]
and the final convolution is
\[
 R(4)=\sum_{j=0}^{4}U_jV_{4-j}
     =1\cdot48+2\cdot27+6\cdot12+9\cdot4+16\cdot1=226.
\]
 before \appendix.
% Uses only the amsart preamble of Diagonal Rational Partitions.
\section{Exact computation by denominator elimination}
\label{sec:exact-computation}

Fix \(N\ge2\) and an integer \(M\ge0\). We compute all counts
\(R_N(m)\), \(0\le m\le M\), in one pass. The last entry for \(N=M=n\)
is \(R(n)\). If the remaining denominators divide \(\Lambda\), their
parts have total in \(\Lambda^{-1}\mathbb Z\). A partial sum that can
be completed to an integer must therefore belong to the same lattice.

Let \(b_1,\ldots,b_d\) be a permutation of \(2,\ldots,N\), where
\(d=N-1\), and put
\[
  \Lambda_d=1,\qquad
  \Lambda_{i-1}=\operatorname{lcm}(b_i,\Lambda_i)
  \quad (i=d,d-1,\ldots,1).
\]
Thus \(\Lambda_i\) is the least common multiple of the denominators
following \(b_i\). We order the denominators by decreasing greatest
prime factor, so that after the group for \(p\) all remaining
denominators are coprime to \(p\).

\subsection{The recurrence}
Let \(f_j\) count the retained configurations of total \(j/D\).
To process denominator \(b\), refine the lattice to denominator
\(H=\operatorname{lcm}(D,b)\). Each allowed part \(a/b\) now has
integer weight \(s=aH/b\). After all parts with denominator \(b_i\)
have been included, the partial sums eligible for retention lie in
\begin{equation}
  H^{-1}\mathbb Z\cap\Lambda_i^{-1}\mathbb Z
  =\gcd(H,\Lambda_i)^{-1}\mathbb Z.
  \label{eq:exact-lattice-intersection}
\end{equation}
The projection therefore keeps every \(H/\gcd(H,\Lambda_i)\)-th
coefficient, starting at index zero.

\smallskip
\noindent\begin{minipage}{\linewidth}
\textsc{RationalPartitionRow}\((N,M)\)
\begin{quote}
\small
\begin{tabbing}
\hspace{1.2em}\=\hspace{1.2em}\=\hspace{1.2em}\=\kill
Choose \(b_1,\ldots,b_d\) and compute \(\Lambda_d,\ldots,\Lambda_0\).\\
\(D\gets1\), \quad \(f\gets(1,0,\ldots,0)\) of length \(M+1\).\\
\textbf{for} \(i=1,\ldots,d\) \textbf{do}\\
\> \(b\gets b_i\), \quad \(H\gets\operatorname{lcm}(D,b)\).\\
\> Set \(g_0,\ldots,g_{MH}\) to zero.\\
\> \textbf{for} \(j=0,\ldots,MD\) \textbf{do} \(g_{jH/D}\gets f_j\).\\
\> \textbf{for} \(1\le a\le\min(N,Mb)\) with \(\gcd(a,b)=1\) \textbf{do}\\
\>\> \(s\gets aH/b\).\\
\>\> \textbf{for} \(j=s,\ldots,MH\) in increasing order \textbf{do}\\
\>\>\> \(g_j\gets g_j+g_{j-s}\).\\
\> \(D\gets\gcd(H,\Lambda_i)\), \quad \(h\gets H/D\).\\
\> \(f\gets(g_{kh})_{k=0}^{MD}\).\\
\textbf{return} \(f\).
\end{tabbing}
\end{quote}
\end{minipage}

\begin{proposition}\label{prop:exact-projection}
The algorithm returns \((R_N(m))_{m=0}^{M}\) for every ordering of the
denominators.
\end{proposition}
\begin{proof}
The initial array counts the empty multiset, and refinement carries
each total \(j/D\) to index \(jH/D\). For a new part of weight \(s\),
let \(u_j,v_j\) be the counts before and after its inclusion. Removing
one copy from each multiset that contains the new part gives
\[
  v_j=u_j+v_{j-s},\qquad v_k=0\quad(k<0).
\]
The increasing loop implements this recurrence. Each multiset is
counted once, since every part is introduced once.

Fix a partition of \(m\le M\), with subtotal \(x\) after \(b_i\).
Since \(m-x\in\Lambda_i^{-1}\mathbb Z\) and \(m\) is an integer,
\(x\in\Lambda_i^{-1}\mathbb Z\). If its previous subtotal was retained,
the updates reach \(x\) on the lattice \(H^{-1}\mathbb Z\).
By \eqref{eq:exact-lattice-intersection}, \(Hx\) is a multiple of
\(h\), so this subtotal is retained too. Positivity gives \(x\le m\),
which justifies truncation at \(M\). By induction every partition of
\(m\) survives. Finally \(\Lambda_d=1\), hence \(D=1\) and \(f_m=R_N(m)\).
\end{proof}

\subsection{The count at order three}
Take the denominators in the order \(3,2\). The thirds must have an
integer subtotal, since \(\gcd(3,2)=1\). To obtain a subtotal \(j\)
from \(1/3,2/3\), their multiplicities satisfy \(u+2v=3j\). For
\(1/2,3/2\) the corresponding equation is \(u+3v=2j\). Thus the two
counts are \(T_j=\lfloor3j/2\rfloor+1\) and
\(B_{2j}\), where \(B_s=\lfloor s/3\rfloor+1\) counts a subtotal
\(s/2\) from the half parts. For \(j=0,1,2,3\) this gives
\[
 (T_j)=(1,2,4,5),\qquad (B_{2j})=(1,1,2,3).
\]
Pairing complementary subtotals gives
\[
 R(3)=\sum_{j=0}^{3}T_jB_{2(3-j)}
     =1\cdot3+2\cdot2+4\cdot1+5\cdot1=16.
\]

\subsection{The count at order four}
Use the order \(3,4,2\). The thirds again have an integer subtotal,
since the remaining denominators divide \(4\). There is now a third
part with denominator \(3\), namely \(4/3\). Its inclusion gives the
count
\[
 U_j=[z^{3j}]\frac{1}{(1-z)(1-z^2)(1-z^4)}.
\]
For the quarter parts \(1/4,3/4\), the next denominator is \(2\),
so \eqref{eq:exact-lattice-intersection} retains half-integer
subtotals. For instance, two copies of \(1/4\) can be completed to
\(1\) by one copy of \(1/2\). The number of ways to make \(r/2\)
from \(1/4,3/4\) is
\[
 Q_r=\#\{(u,v)\in\mathbb Z_{\ge0}^2:u+3v=2r\}
     =\left\lfloor\frac{2r}{3}\right\rfloor+1.
\]
The half parts are again \(1/2,3/2\). Thus the combined count for
integer subtotal \(j\) is \(V_j=\sum_{r=0}^{2j}Q_r B_{2j-r}\), giving
\[
 \begin{array}{c|rrrrr}
 j   &0&1&2&3&4\\ \hline
 U_j &1&2&6&9&16\\
 V_j &1&4&12&27&48
 \end{array}
\]
and the final convolution is
\[
 R(4)=\sum_{j=0}^{4}U_jV_{4-j}
     =1\cdot48+2\cdot27+6\cdot12+9\cdot4+16\cdot1=226.
\]
 before \appendix.
% Uses only the amsart preamble of Diagonal Rational Partitions.
\section{Exact computation by denominator elimination}
\label{sec:exact-computation}

Fix \(N\ge2\) and an integer \(M\ge0\). We compute all counts
\(R_N(m)\), \(0\le m\le M\), in one pass. The last entry for \(N=M=n\)
is \(R(n)\). If the remaining denominators divide \(\Lambda\), their
parts have total in \(\Lambda^{-1}\mathbb Z\). A partial sum that can
be completed to an integer must therefore belong to the same lattice.

Let \(b_1,\ldots,b_d\) be a permutation of \(2,\ldots,N\), where
\(d=N-1\), and put
\[
  \Lambda_d=1,\qquad
  \Lambda_{i-1}=\operatorname{lcm}(b_i,\Lambda_i)
  \quad (i=d,d-1,\ldots,1).
\]
Thus \(\Lambda_i\) is the least common multiple of the denominators
following \(b_i\). We order the denominators by decreasing greatest
prime factor, so that after the group for \(p\) all remaining
denominators are coprime to \(p\).

\subsection{The recurrence}
Let \(f_j\) count the retained configurations of total \(j/D\).
To process denominator \(b\), refine the lattice to denominator
\(H=\operatorname{lcm}(D,b)\). Each allowed part \(a/b\) now has
integer weight \(s=aH/b\). After all parts with denominator \(b_i\)
have been included, the partial sums eligible for retention lie in
\begin{equation}
  H^{-1}\mathbb Z\cap\Lambda_i^{-1}\mathbb Z
  =\gcd(H,\Lambda_i)^{-1}\mathbb Z.
  \label{eq:exact-lattice-intersection}
\end{equation}
The projection therefore keeps every \(H/\gcd(H,\Lambda_i)\)-th
coefficient, starting at index zero.

\smallskip
\noindent\begin{minipage}{\linewidth}
\textsc{RationalPartitionRow}\((N,M)\)
\begin{quote}
\small
\begin{tabbing}
\hspace{1.2em}\=\hspace{1.2em}\=\hspace{1.2em}\=\kill
Choose \(b_1,\ldots,b_d\) and compute \(\Lambda_d,\ldots,\Lambda_0\).\\
\(D\gets1\), \quad \(f\gets(1,0,\ldots,0)\) of length \(M+1\).\\
\textbf{for} \(i=1,\ldots,d\) \textbf{do}\\
\> \(b\gets b_i\), \quad \(H\gets\operatorname{lcm}(D,b)\).\\
\> Set \(g_0,\ldots,g_{MH}\) to zero.\\
\> \textbf{for} \(j=0,\ldots,MD\) \textbf{do} \(g_{jH/D}\gets f_j\).\\
\> \textbf{for} \(1\le a\le\min(N,Mb)\) with \(\gcd(a,b)=1\) \textbf{do}\\
\>\> \(s\gets aH/b\).\\
\>\> \textbf{for} \(j=s,\ldots,MH\) in increasing order \textbf{do}\\
\>\>\> \(g_j\gets g_j+g_{j-s}\).\\
\> \(D\gets\gcd(H,\Lambda_i)\), \quad \(h\gets H/D\).\\
\> \(f\gets(g_{kh})_{k=0}^{MD}\).\\
\textbf{return} \(f\).
\end{tabbing}
\end{quote}
\end{minipage}

\begin{proposition}\label{prop:exact-projection}
The algorithm returns \((R_N(m))_{m=0}^{M}\) for every ordering of the
denominators.
\end{proposition}
\begin{proof}
The initial array counts the empty multiset, and refinement carries
each total \(j/D\) to index \(jH/D\). For a new part of weight \(s\),
let \(u_j,v_j\) be the counts before and after its inclusion. Removing
one copy from each multiset that contains the new part gives
\[
  v_j=u_j+v_{j-s},\qquad v_k=0\quad(k<0).
\]
The increasing loop implements this recurrence. Each multiset is
counted once, since every part is introduced once.

Fix a partition of \(m\le M\), with subtotal \(x\) after \(b_i\).
Since \(m-x\in\Lambda_i^{-1}\mathbb Z\) and \(m\) is an integer,
\(x\in\Lambda_i^{-1}\mathbb Z\). If its previous subtotal was retained,
the updates reach \(x\) on the lattice \(H^{-1}\mathbb Z\).
By \eqref{eq:exact-lattice-intersection}, \(Hx\) is a multiple of
\(h\), so this subtotal is retained too. Positivity gives \(x\le m\),
which justifies truncation at \(M\). By induction every partition of
\(m\) survives. Finally \(\Lambda_d=1\), hence \(D=1\) and \(f_m=R_N(m)\).
\end{proof}

\subsection{The count at order three}
Take the denominators in the order \(3,2\). The thirds must have an
integer subtotal, since \(\gcd(3,2)=1\). To obtain a subtotal \(j\)
from \(1/3,2/3\), their multiplicities satisfy \(u+2v=3j\). For
\(1/2,3/2\) the corresponding equation is \(u+3v=2j\). Thus the two
counts are \(T_j=\lfloor3j/2\rfloor+1\) and
\(B_{2j}\), where \(B_s=\lfloor s/3\rfloor+1\) counts a subtotal
\(s/2\) from the half parts. For \(j=0,1,2,3\) this gives
\[
 (T_j)=(1,2,4,5),\qquad (B_{2j})=(1,1,2,3).
\]
Pairing complementary subtotals gives
\[
 R(3)=\sum_{j=0}^{3}T_jB_{2(3-j)}
     =1\cdot3+2\cdot2+4\cdot1+5\cdot1=16.
\]

\subsection{The count at order four}
Use the order \(3,4,2\). The thirds again have an integer subtotal,
since the remaining denominators divide \(4\). There is now a third
part with denominator \(3\), namely \(4/3\). Its inclusion gives the
count
\[
 U_j=[z^{3j}]\frac{1}{(1-z)(1-z^2)(1-z^4)}.
\]
For the quarter parts \(1/4,3/4\), the next denominator is \(2\),
so \eqref{eq:exact-lattice-intersection} retains half-integer
subtotals. For instance, two copies of \(1/4\) can be completed to
\(1\) by one copy of \(1/2\). The number of ways to make \(r/2\)
from \(1/4,3/4\) is
\[
 Q_r=\#\{(u,v)\in\mathbb Z_{\ge0}^2:u+3v=2r\}
     =\left\lfloor\frac{2r}{3}\right\rfloor+1.
\]
The half parts are again \(1/2,3/2\). Thus the combined count for
integer subtotal \(j\) is \(V_j=\sum_{r=0}^{2j}Q_r B_{2j-r}\), giving
\[
 \begin{array}{c|rrrrr}
 j   &0&1&2&3&4\\ \hline
 U_j &1&2&6&9&16\\
 V_j &1&4&12&27&48
 \end{array}
\]
and the final convolution is
\[
 R(4)=\sum_{j=0}^{4}U_jV_{4-j}
     =1\cdot48+2\cdot27+6\cdot12+9\cdot4+16\cdot1=226.
\]
